%%% SDD - Organisation: Beginn
\begin{center}
  \makeatletter
  \textbf{
    \@title
  }
  \makeatother
  \\[1em]
  Informatik -- Eingebettete Systeme -- Hochschule Bremerhaven
\end{center}

\begin{table}[ht]
  \centering
  \begin{tabular}[t]{|p{15mm}|p{30mm}|p{30mm}|p{50mm}|r|}
    \hline
    \textbf{Version} & \textbf{Autor*inn*en} & \textbf{geändert} & \textbf{Beschreibung} & \textbf{fertig am}\\
    \hline
    0.1              & A. B.            &  Abschnitte 8 bis 12  & Erstellung des Dokumentes & 01.10.2019\\
    \hline
  \end{tabular}
  \caption{Versionshistorie}
\end{table}

\vspace*{4em}

Relevante Standards und Normen
\begin{itemize}
\item IEEE 1016-2009 -- IEEE Standard for Information Technology --
  Systems Design -- Software Design Descriptions

  \cite{ieee-1016-2009-software-design-description}
\end{itemize}
\clearpage
%%% SDD - Organisation: Ende

\chapter{Einleitung}
\label{sdd:einleitung}

\section{Designübersicht}
\label{sdd:designuebersicht}

\begin{figure}[htp]
  \centering

  \textbf{\hinweis{Hierhin gehört die Bausteinsicht, die die einzelnen Komponenten auf
  der höchsten Ebene zeigt.}}

  \caption{Bausteinsicht}
  \label{sdd:bausteinsicht}
\end{figure}

\section{Anforderungsverfolgungsmatrix}
\label{sdd:anforderungsmatrix}

\begin{tabular}[t]{|l||l|l|l|l|l|}
  Anforderung & \hinweis{Baustein-1} & \hinweis{Baustein-2} & \hinweis{Baustein-3} & \hinweis{Baustein-4} & \dots\\
  \hline
  \hline
  \hinweis{REQ-EXT-USER-1} &         &            &            &            & \\
  \hinweis{REQ-EXT-USER-2} &         &            &            &            & \\
  \hline
  \hinweis{REQ-EXT-HW-1}   &         &            &            &            & \\
  \hinweis{REQ-EXT-HW-2}   &         &            &            &            & \\
  \hinweis{REQ-EXT-HW-3}   &         &            &            &            & \\
  \hline
  \hinweis{REQ-EXT-PROD-1} &         &            &            &            & \\
  \dots          &         &            &            &            & \\
  \hline
\end{tabular}

\chapter{Systemarchitektur}
\label{sdd:systemarchitektur}

\section{Gewählte Systemarchitektur}
\label{sdd:systemarchitektur-gewaehlt}
\hinweis{Hier steht die Architektur, die Sie ausgewählt haben, wenn es
  andere Möglichkeiten gibt, die Sie betrachtet haben, dann stehen
  diese im nächsten Punkt, der Diskussion alternativer Designs}

\section{Diskussion alternativer Designs}
\label{sdd:systemarchitektur-alternativen}
\hinweis{Dieser Abschnitt enthält alles, was Sie zur Lösungsfindung
  probiert und diskutiert haben. Insbesondere schreiben Sie hier auch
  Experimente, deren Aufbau, Durchführung und deren Ergebnisse auf.}

\subsection{Experiment-N}
\label{sec:experimente-N}

\subsubsection{Zweck}
\label{sec:experimente-zweck-N}

\subsubsection{Aufbau}
\label{sec:experimente-aufbau-N}

\subsubsection{Durchführung}
\label{sec:experimente-durchfuehrung-N}

\subsubsection{Ergebisse}
\label{sec:experimente-ergebisse-N}


\section{Systemschnittstellenbeschreibung}
\label{sdd:systemschnittstellen}

% Kopiere Schnittstellenbeschreibung von Beginn bis Ende für jede Schnittstelle
% Schnittstellenbeschreibung: Beginn
\subsection{Schnittstelle-N}
\label{sdd:schnittstellen-N}
\hinweis{Gliedern Sie die Schnittstellenbeschreibung passend auf.}

\subsubsection{Kommunikation: Initiierung, Nachrichten und Ende}
\label{sec:schnittstellen-initiierung-N}
\hinweis{Hier steht, welche Komponente, die Kommunikation startet und
  beendet und welches die Nachrichten sind, die über die Schnittstelle
  ausgetauscht werden. Sequenzdiagramme gehören hier hin.}

\subsubsection{Zeitverhalten}
\label{sec:schnittstellen-zeitverhalten-N}
\hinweis{Hier steht, wie sich die Beteiligten der Schnittstelle
  zeitlich verhalten sollen. Z.B.
  \begin{itemize}
  \item bei sehr vielen Nachrichten (\enquote{Meldungsschauer}
    (engl. \enquote{burst}))
  \item Timeouts
  \item Timing-Diagramme!
  \end{itemize}
}

\subsubsection{Fehlerbehandlung}
\label{sec:schnittstellen-fehlerbehandlung-N}


% Schnittstellenbeschreibung: Ende

\chapter{Benutzerschnittstellendesign}
\label{sdd:benutzerschnittstellendesign}
\hinweis{Benutzerschnittstellen sind durchaus auch
  Systemschnittstellen. Sie sind besonders, da hier nicht mit
  technischen Systemen interagiert wird, sondern mit Menschen. Man
  könnte diesen Punkt auch Mensch/Maschine-Kommunikation bezeichnen.
  Taster und LEDs gehören zur Benutzerschnittstelle.}

\section{Beschreibung der Benutzerschnittstelle}
\label{sdd:benutzerschnittellenbeschreibung}

\subsection{Bildschirmfotos / Fotos der Benutzeroberfläche}
\label{sdd:benutzeroberfläche}

\subsection{Objekte und Aktionen}
\label{sdd:benutzerschnittstelle-objekte-aktionen}


\chapter{Detaillierte Beschreibung der Komponenten}
\label{sdd:detailbeschreibung}

% Kopiere Komponentenbeschreibung von Beginn bis Ende für jede Komponente
% Komponentenbeschreibung: Beginn
\section{Baustein/Komponente-N}
\label{sdd:komponenten-N}
\hinweis{
  Hier informieren Sie über
  \begin{itemize}
  \item die interne Struktur des Bausteins/der Komponente, in Hardware
    und in Software
  \item die physikalischen und logischen Eigenschaften und deren
    Aufteilung in Hardware und Software
  \item interne Zustände (sofern es das gibt)
  \item interne Abläufe (Prozesse)
  \item Datenstrukturen
  \item Initialisierung
  \item \dots
  \end{itemize}
}

% Komponentenbeschreibung: Ende

\chapter{Zusatzmaterial}
\label{sdd:zusatzmaterial}

